% ===================================================================
\section{Diapositivas de respaldo}

Están después de la diapositiva de cierre y tienen numeración propia (1/6 a 6/6). No
cuentan en el tiempo. Úselas solo para responder preguntas.

\subsection*{Respaldo 1 · Modelo de la cartera multicapa}
\textbf{Cuándo:} si preguntan cómo es exactamente el modelo de la tarea 6 o de dónde
salen los parámetros.

\textbf{Qué decir:} el modelo es el mismo de la diapositiva 5 con tres cambios. El tiempo
está en meses y el déficit es una fracción del ingreso total del cliente. El contagio y
el acople tienen un valor por capa: la nómina transmite más que las transferencias. Y
hay ruido aleatorio e impagos propios de cada cliente según su probabilidad de impago.
El colchón es más pequeño para los clientes más riesgosos. Los valores son ilustrativos y
están en el archivo \texttt{configs/credit.yaml}.

\begin{terminos}
\begin{description}
\item[$\ell$.] Capa: N (nómina), S (proveedores), T (transferencias).
\item[$W^\ell_{ij}$.] Flujo mensual de $j$ a $i$ en la capa $\ell$, en salarios mínimos.
\item[$s_i$.] Ingreso total del cliente, incluido el que llega de fuera del banco.
\item[$k_i$.] Ingreso que llega por la red del banco.
\item[$\beta_\ell$, $K_\ell$.] Contagio y acople de fase de cada capa.
\item[$\sigma\xi_i(t)$.] Ruido aleatorio con intensidad $\sigma$. Se integra con el método
de Euler–Maruyama.
\item[$x_i^{(c)}$.] Colchón $c\,(\mathrm{PD}_\text{ref}/\mathrm{PD}_i)^\kappa\varepsilon_i$:
más pequeño cuanto mayor es la probabilidad de impago $\mathrm{PD}_i$.
\item[$\lambda$.] Déficit que genera perder a un pagador que aporta todo el ingreso: un
mes de ingreso.
\item[Impago propio.] Cada cliente tiene una fecha de impago aleatoria con tasa
$-\ln(1-\mathrm{PD}_i)/12$ por mes, de modo que cae en el año con probabilidad
$\mathrm{PD}_i$.
\end{description}
\end{terminos}

\subsection*{Respaldo 2 · Métricas completas}
\textbf{Cuándo:} si preguntan por otras métricas o por el mundo nulo.

\textbf{Qué decir:} todas las métricas mejoran en la misma dirección. La PR-AUC, que es
la más exigente cuando los impagos son pocos, pasa de 0.277 a 0.323. En el mundo nulo los
niveles son iguales. El impago a doce meses es de 9.6~\% con red y de 5.9~\% en el nulo.

\begin{terminos}
\begin{description}
\item[PR-AUC.] Área bajo la curva de precisión contra exhaustividad. Un modelo al azar
obtiene la tasa de impago, cerca de 0.1.
\item[KS.] Estadístico de Kolmogorov–Smirnov: máxima separación entre las distribuciones
del score de los que caen y de los que no. Muy usado en banca.
\item[Brier.] Error cuadrático medio de la probabilidad predicha. Menor es mejor.
\end{description}
\end{terminos}

\subsection*{Respaldo 3 · Calibración y relación sincronía–impago}
\textbf{Cuándo:} si preguntan si las probabilidades son confiables, o si basta con mirar
la coherencia.

\textbf{Qué decir:} a la izquierda, ambos modelos siguen la diagonal: cuando predicen
10~\%, cae cerca del 10~\%. El híbrido separa mejor a los más riesgosos: en su decil más
alto cae el 35~\%, frente al 30~\% del clásico. A la derecha, con red, el impago pasa de
4.4~\% en los bien sincronizados a entre 11 y 14~\% en los desincronizados. En el nulo
también hay diferencia, pero menor, porque los informales son más riesgosos y están menos
sincronizados. Por eso la prueba es la comparación de modelos y no esta gráfica.

\begin{terminos}
\begin{description}
\item[Calibración.] Coincidencia entre la probabilidad predicha y la frecuencia
observada.
\item[Decil.] Cada décima parte de los clientes, ordenados por riesgo.
\item[Variable de confusión.] Variable que influye a la vez en la causa aparente y en el
efecto; aquí, la informalidad.
\end{description}
\end{terminos}

\subsection*{Respaldo 4 · Redistribución y recuperación (tarea 2)}
\textbf{Cuándo:} si preguntan por la tarea 2, que la exposición solo menciona.

\textbf{Qué decir:} el Laplaciano traslada déficit sin crearlo; lo verificamos con
precisión de $10^{-10}$. La estabilidad la da $\sigma$, el mayor valor propio del sistema
completo, y el tiempo de recuperación es $-1/\sigma$. A la izquierda, las redes mejor
conectadas se recuperan más rápido. A la derecha, comunidades encontradas con el vector de
Fiedler.

\begin{terminos}
\begin{description}
\item[$\sigma$.] Mayor parte real de los valores propios de $M=-\alpha I-D_0L+\beta W$.
Negativa: la red se recupera.
\item[$\tau_\text{rec}$.] Tiempo de recuperación, $-1/\sigma$.
\item[$\lambda_2$.] Conectividad algebraica: segundo valor propio más pequeño del
Laplaciano de la red simetrizada; vale cero si la red está partida.
\item[Vector de Fiedler.] Vector propio asociado a $\lambda_2$; su signo divide la red en
dos grupos poco conectados.
\end{description}
\end{terminos}

\subsection*{Respaldo 5 · Tarea 5 en detalle}
\textbf{Cuándo:} si preguntan por redes neuronales o por cómo se generaron los datos.

\textbf{Qué decir:} las simulaciones cubren 13 parámetros con un hipercubo latino, con
más puntos cerca de $R=1$, donde el comportamiento cambia. Para la cascada de cada
entidad, una red neuronal sobre grafos explica el 60~\% de la variación. La PINN estima
los parámetros del modelo a partir de trayectorias: recupera $\gamma$ y $K$ con menos de
1~\% de error; $\beta$ tiene 23~\% porque su efecto se confunde con el de $\alpha$.

\begin{terminos}
\begin{description}
\item[Hipercubo latino.] Muestreo que divide el rango de cada parámetro en intervalos y
usa cada intervalo una sola vez.
\item[GCN.] \emph{Graph Convolutional Network}: red neuronal que combina la información
de cada nodo con la de sus vecinos (Kipf y Welling, 2017).
\item[PINN.] \emph{Physics-Informed Neural Network}: red neuronal entrenada para ajustar
los datos y cumplir las ecuaciones del modelo (Raissi et al., 2019).
\item[$R^2$.] Fracción de la variación explicada por el modelo.
\item[$\rho_s$.] Correlación de Spearman: compara órdenes, no valores.
\end{description}
\end{terminos}

\subsection*{Respaldo 6 · Referencias principales}
\textbf{Cuándo:} si piden las fuentes. La bibliografía completa está en
\texttt{docs/bibliografia.md}.

% ===================================================================
\section{Preguntas generales}

\begin{preguntas}
\pr{¿Cuál es el aporte original?}{Mostrar con un experimento controlado que la
sincronización de los ciclos de pago, estimada solo con transacciones, contiene
información sobre el impago que no está ni en el score ni en la red estática. Incluye el
método para estimarla y un control que descarta artefactos.}
\pr{¿Por qué un enfoque desde la física?}{Porque las herramientas son de la física:
sistemas dinámicos, valores propios, osciladores acoplados, transiciones y avalanchas.
Permiten responder cuándo la red amplifica y cómo se sincroniza.}
\pr{¿El resultado es causal?}{En la simulación sabemos la causa porque la pusimos, y el
experimento de $\gamma$ muestra que la señal depende de ella. Con datos reales el método
sería predictivo; probar causalidad requeriría otro diseño.}
\pr{¿Cómo se evitó la fuga de información?}{Las variables solo usan pagos de los primeros
seis meses, y el código rechaza pagos posteriores. La probabilidad verdadera está
prohibida para los modelos. Las carteras de prueba nunca se usan para entrenar. Y el
mundo nulo da mejora cero.}
\pr{¿Qué pasa con otros parámetros?}{El tamaño de la mejora cambia: crece con $\gamma$ y
desaparece sin él. El patrón se mantiene. El experimento completo se repite en unos 20
segundos.}
\pr{¿Y si el banco no ve toda la red?}{No se probó. Es trabajo futuro: un banco ve solo a
sus clientes, mientras que el operador de Bre-B vería todos los pagos.}
\pr{¿Por qué Kuramoto?}{Es el modelo mínimo de sincronización con resultados exactos
conocidos, lo que permite validar el código, y su interpretación es directa: cada
entidad ajusta su ciclo al de quienes le pagan.}
\end{preguntas}

% ===================================================================
\section{Glosario}

\begin{longtable}{p{0.22\textwidth}p{0.72\textwidth}}
\toprule
\textbf{Término} & \textbf{Significado} \\
\midrule
\endhead
AUC & Probabilidad de que el modelo asigne más riesgo a quien cae que a quien no cae.
0.5 es azar; 1 es perfecto. \\
Bre-B & Sistema de pagos inmediatos del Banco de la República de Colombia. \\
Cascada & Impagos en cadena. \\
Coherencia ($r_i$) & Qué tan a tiempo llega el ingreso de cada entidad respecto a su
propio ciclo; de $-1$ a $1$. \\
Colchón & Déficit máximo que una entidad soporta antes de dejar de pagar. \\
Control nulo & Mundo idéntico sin el mecanismo; mide los artefactos del método. \\
DANE & Departamento Administrativo Nacional de Estadística. \\
DebtRank & Medida de propagación de estrés financiero. \\
Déficit ($x_i$) & Liquidez que le falta a una entidad. \\
Desfase & Diferencia entre la llegada del ingreso y el vencimiento de las
obligaciones. \\
Fase ($\theta_i$) & Punto del ciclo de pagos en que está una entidad. \\
GCN & Red neuronal sobre grafos. \\
Gini & $2\,\mathrm{AUC}-1$; métrica habitual en banca. \\
GroupKFold & Validación cruzada que mantiene cada grupo (red o cartera) en una sola
partición. \\
Intervalo de confianza & Rango que contiene el valor verdadero con la confianza
indicada (aquí 95~\%). \\
Kuramoto & Modelo de osciladores acoplados para estudiar la sincronización. \\
Laplaciano & Matriz que describe flujos que se conservan en una red. \\
ML & Aprendizaje automático (\emph{machine learning}). \\
PD & Probabilidad de impago a un año. \\
PINN & Red neuronal informada por la física. \\
$R$ & Índice de amplificación $\beta\rho(W)/\alpha$; menor que 1, la red amortigua. \\
Radio espectral & Mayor valor propio de una matriz en módulo. \\
Score crediticio & Modelo que estima la probabilidad de impago de un cliente. \\
SHAP & Método que explica cuánto aporta cada variable a una predicción. \\
SMMLV & Salario mínimo mensual legal vigente. \\
XGBoost & Conjunto de árboles de decisión entrenados en secuencia. \\
\bottomrule
\end{longtable}
